\chapter{Concurrency, Operating Systems and Networks}\label{ch:iv:concurrency-operating-systems-and-networks}

Two threads each add one to a shared counter a million times; the program prints a number well below two million. The
candidate is asked for the smallest value it could print, which is not what anyone expects, and then for three ways to make it
print two million, ranked by what each costs on a hot path. Systems questions at trading firms follow the path of a message
through a machine: threads and the memory they share, the kernel and its system calls, virtual memory and caches, and the
network stack that delivers the packet. The engineering is One Quant Book~13's (chapters 3, 11 to~13) and Book~14's; this
chapter is the interview.

\section{Atomics, locks and memory ordering}

Two threads that access the same memory, at least one writing, without synchronisation, have a data race; in C++ and Rust that
is undefined behaviour (One Quant Book~13, chapter~9), and ThreadSanitizer reports it. The classic lost update comes from
\texttt{++counter} being a load, an add and a store: two threads can load the same value and both store its successor. Under the
simplest model of an interleaving (each thread's steps in order, one global order of memory operations: sequential
consistency) the outcomes can be enumerated exhaustively, which is the interview's method for small programs.

\begin{proposition}[How low can the counter go?]\label{prop:iv:concurrency-operating-systems-and-networks:counter}
If two threads each perform $n \ge 2$ increments by a separate load and store, the final value under sequential consistency is
at least 2 and at most $2n$, and 2 is reached. For $n = 1$ it is 1 or 2.
\end{proposition}

\begin{proof}
Every store writes a loaded value plus one, so once a thread has stored, the counter is at least 1 forever. The final value is the
last store, which is some thread's last increment; for $n \ge 2$ that increment's load comes after the same thread's first store, so
it reads at least 1 and stores at least 2. To reach 2: thread A loads 0; thread B runs $n - 1$ increments; A stores 1; B loads 1; A
runs its remaining $n - 1$ increments; B stores 2. The chapter's code enumerates every interleaving for $n \le 4$ and finds every
value from 2 to $2n$.
\end{proof}

Fixes, from cheapest to dearest on a hot path: give each thread its own counter and sum at the end (no sharing at all); an atomic
\texttt{fetch\_add} (a locked instruction on one cache line that bounces between cores under contention); a mutex around the
increment (a system call when contended). Memory ordering (One Quant Book~13, chapter~11) decides what else becomes visible with an
atomic operation: \texttt{relaxed} guarantees only atomicity, \texttt{release}/\texttt{acquire} make writes before a release visible
after the matching acquire (the message-passing pattern), and \texttt{seq\_cst} adds a single total order, needed for patterns
such as store buffering, where each thread writes one flag and reads the other's.

\omcode{code/interviews/25-concurrency-operating-systems-and-networks/cpp/iv_spsc.hpp}{10}{33}{A single-producer
single-consumer queue: the release store of the index publishes the element; the two indices live on separate cache
lines.}\label{lst:iv:concurrency-operating-systems-and-networks:spsc}

\section{Processes, threads, scheduling and system calls}

A system call crosses from user to kernel mode and back; a context switch saves one thread's registers and restores another's,
and pollutes caches and translation tables. Both cost far more than a function call, which is why latency-critical threads are
pinned to isolated cores, poll instead of sleeping (busy polling), and avoid the kernel on the data path (kernel bypass, One Quant
Book~13, chapter~13). A mutex that is uncontended is a couple of atomic instructions; contended, it sleeps in the kernel; a spinlock
never sleeps and suits very short critical sections on dedicated cores only.

\section{Memory: virtual memory, page faults, caches and false sharing}

Virtual addresses are translated through a radix tree of page tables: with 48-bit addresses, 4 KiB pages and 8-byte entries, each
level resolves 9 bits and four levels plus a 12-bit offset cover the address. The translation lookaside buffer caches
translations; its reach is its entry count times the page size, so huge pages (2 MiB) multiply it by 512. The first touch of a page
faults (One Quant Book~13, chapter~6), which is why trading processes pre-fault their memory at start-up. Caches move 64-byte lines;
two threads writing different variables in the same line make it bounce between cores (false sharing), which padding each hot
variable to its own line removes.

\section{Networks: what happens when a packet arrives}

\begin{omfigure}
\begin{tikzpicture}[font=\footnotesize,>={Stealth},
  box/.style={draw=omDef,rounded corners=2pt,fill=omDef!6,minimum height=0.75cm,text width=2.35cm,align=center,inner sep=2pt},
  byp/.style={draw=omThm,rounded corners=2pt,fill=omThm!8,minimum height=0.75cm,text width=2.35cm,align=center,inner sep=2pt}]
\node[box] (nic) at (0,0) {NIC receives frame, DMA into a ring};
\node[box] (irq) at (3.1,1.3) {interrupt, then NAPI polling};
\node[box] (stk) at (6.2,1.3) {IP and UDP processing};
\node[box] (sock) at (9.3,1.3) {socket buffer};
\node[box] (app) at (12.7,0) {application};
\node[box] (sys) at (9.3,-0.05) {\texttt{recvmsg}: system call, copy};
\node[byp] (ring) at (6.2,-1.35) {user-space ring, polled by the application};
\draw[->,thick,omDef] (nic) -- (irq);
\draw[->,thick,omDef] (irq) -- (stk);
\draw[->,thick,omDef] (stk) -- (sock);
\draw[->,thick,omDef] (sock) -- (sys);
\draw[->,thick,omDef] (sys) -- (app);
\draw[->,thick,omThm] (nic) -- (ring);
\draw[->,thick,omThm] (ring) -- (app);
\node[font=\scriptsize,omDef,anchor=south] at (6.2,1.8) {kernel path};
\node[font=\scriptsize,omThm,anchor=north] at (6.2,-1.85) {kernel bypass};
\end{tikzpicture}

\omcaption{A received packet's two routes to the application. Through the kernel: the network card writes the frame into memory by
DMA, an interrupt starts the driver's polling, the stack processes the headers and queues the data on the socket, and a system call
copies it out. With kernel bypass the card's ring is mapped into the application, which polls it: no interrupt, no system call, no
copy.}\label{fig:iv:concurrency-operating-systems-and-networks:packet}
\end{omfigure}

TCP gives a reliable ordered byte stream with congestion control and retransmission; UDP gives datagrams that may be lost, duplicated
or reordered, with no head-of-line blocking. Market data is multicast over UDP (One Quant Book~13, chapter~16), because one sender
reaches many receivers and a lost packet should not delay the next; receivers detect gaps with sequence numbers and recover from a
second line or a retransmission service. Order entry uses TCP, because every message must arrive in order and a lost order must not
go unnoticed.

\section{Worked answers}

\begin{example}[Little's law and a ring buffer]\label{ex:iv:concurrency-operating-systems-and-networks:little}
\emph{``A decoding stage receives 200\,000 messages a second and each spends 15 microseconds in it on average. How many are
inside the stage at once? During a burst the feed delivers a million messages a second for a while and the stage still
drains 200\,000; how long can the burst last before a ring buffer of 1\,024 slots overflows?''} Little's law: the average
number in a stable system is the arrival rate times the time in the system, $200\,000 \times 15 \times 10^{-6} = 3$
messages. In the burst the queue grows at $10^6 - 2 \times 10^5 = 800\,000$ messages a second, so 1\,024 slots last
$1\,024/800\,000 = 1.28$ milliseconds. Then the design questions the numbers raise: feeds burst for longer than that at the
open, so either the buffer grows (at 64 bytes a slot, a million slots is 64 megabytes, which no longer fits in cache), or the
stage gets faster, or it sheds work by conflating updates, and the answer should say which the business can accept.
\end{example}

\section{Question bank}

\omcode{code/interviews/25-concurrency-operating-systems-and-networks/cpp/snippets/race_counter.cpp}{1}{15}{Two threads, one
counter.}\label{lst:iv:concurrency-operating-systems-and-networks:race}

\begin{interviewq}[$\star$ \iqroles{developer} \iqfirm{market maker}]\label{iq:iv:concurrency-operating-systems-and-networks:1}
Why does \cref{lst:iv:concurrency-operating-systems-and-networks:race} usually print less than two million? What does the C++
standard say about the program?
\end{interviewq}

\begin{interviewq}[$\star$ \iqroles{developer} \iqfirm{proprietary firm}]\label{iq:iv:concurrency-operating-systems-and-networks:2}
When would you use a spinlock rather than a mutex?
\end{interviewq}

\begin{interviewq}[$\star$ \iqroles{developer} \iqfirm{market maker}]\label{iq:iv:concurrency-operating-systems-and-networks:3}
What is a context switch, what does it cost beyond the switch itself, and how do trading systems avoid them on the hot path?
\end{interviewq}

\begin{interviewq}[$\star$ \iqroles{developer, trader} \iqfirm{any}]\label{iq:iv:concurrency-operating-systems-and-networks:4}
Why do venues usually publish market data by UDP multicast but take orders on TCP sessions?
\end{interviewq}

\begin{interviewq}[$\star\star$ \iqroles{developer} \iqfirm{proprietary firm}]\label{iq:iv:concurrency-operating-systems-and-networks:5}
In the model where each increment is a separate load and store, what is the smallest value the program of
\cref{lst:iv:concurrency-operating-systems-and-networks:race} could print? Describe the interleaving.
\end{interviewq}

\begin{interviewq}[$\star\star$ \iqroles{developer} \iqfirm{market maker}]\label{iq:iv:concurrency-operating-systems-and-networks:6}
Thread 1 writes \texttt{x = 1} then reads \texttt{y}; thread 2 writes \texttt{y = 1} then reads \texttt{x}; both start at 0. Can both
threads read 0? Answer under sequential consistency, on x86, and with C++ relaxed atomics, and say how to forbid it.
\end{interviewq}

\begin{interviewq}[$\star\star$ \iqroles{developer} \iqfirm{proprietary firm}]\label{iq:iv:concurrency-operating-systems-and-networks:7}
A producer writes a message then sets a flag; a consumer waits for the flag then reads the message. Which memory orderings make
this correct, and which outcome do they forbid?
\end{interviewq}

\begin{interviewq}[$\star\star$ \iqroles{developer} \iqfirm{market maker}]\label{iq:iv:concurrency-operating-systems-and-networks:8}
Two threads each update their own \texttt{int64\_t} counter in an array of eight counters, and the program is slower than with one
thread. Why, and what is the fix?
\end{interviewq}

\begin{interviewq}[$\star\star$ \iqroles{developer} \iqfirm{proprietary firm}]\label{iq:iv:concurrency-operating-systems-and-networks:9}
With 48-bit virtual addresses, 4 KiB pages and 8-byte page-table entries, how many levels of page table are there? A TLB has 1\,536
entries: how much memory does it cover using 4 KiB pages, and using 2 MiB huge pages?
\end{interviewq}

\begin{interviewq}[$\star\star\star$ \iqroles{developer} \iqfirm{market maker}]\label{iq:iv:concurrency-operating-systems-and-networks:10}
Write a lock-free single-producer single-consumer queue and justify every memory ordering in it.
\end{interviewq}

\begin{interviewq}[$\star\star\star$ \iqroles{developer} \iqfirm{proprietary firm}]\label{iq:iv:concurrency-operating-systems-and-networks:11}
What happens, step by step, between a UDP packet arriving at the network card and the application having its payload? Which steps
does kernel bypass remove?
\end{interviewq}

\begin{interviewq}[$\star\star\star$ \iqroles{developer, trader} \iqfirm{crypto firm}]\label{iq:iv:concurrency-operating-systems-and-networks:12}
A multicast feed carries sequence numbers. Write the receiver's gap detector, allowing for late and duplicated packets, and say
what the receiver does when a gap stays open.
\end{interviewq}

\begin{interviewq}[$\star\star\star$ \iqroles{developer} \iqfirm{any}]\label{iq:iv:concurrency-operating-systems-and-networks:13}
Thread A locks the order book then the risk state; thread B locks the risk state then the order book. What can happen, and give
two ways to prevent it.
\end{interviewq}

\begin{omsources}
\item M.~Herlihy and N.~Shavit, \emph{The Art of Multiprocessor Programming}, revised reprint, 2012.
\item J.~Postel, RFC 768, \emph{User Datagram Protocol}, 1980; W.~Eddy (ed.), RFC 9293, \emph{Transmission Control Protocol (TCP)},
2022.
\item ISO/IEC 14882:2020, [intro.races] and [atomics.order]; One Quant Book~13, chapters 3, 6, 9, 11--13 and~16.
\end{omsources}
